\subsection{Pregunta 5}

Utilizaron el Grupo 5 de Diffie-Hellman. Investiga por qué hoy en día organizaciones como la NSA o el NIST recomiendan dejar de usar los grupos 2 y 5, y qué grupos (números) se consideran seguros para estándares de cifrado modernos.

\begin{respuesta}
El grupo de Diffie-Hellman define qué tan grandes son los números que se usan
para generar la llave, y de eso depende qué tan difícil es romperla. Los grupos
2 y 5 ya se quedaron cortos:

\begin{itemize}
    \item El \textbf{grupo 2} usa números de 1024 bits y el \textbf{grupo 5} de
    1536 bits. Antes alcanzaba, pero hoy las computadoras son mucho más
    potentes y los ataques son mejores.
    \item En 2015 salió el ataque \textbf{Logjam}, que demostró que los grupos
    de 1024 bits se pueden romper si alguien tiene mucho poder de cómputo (como
    un gobierno). Como además mucha gente usa los mismos números primos
    estándar, el atacante hace el trabajo pesado una sola vez y lo reutiliza
    para espiar un montón de conexiones.
    \item Por eso el NIST y la NSA recomiendan dejarlos y usar grupos más
    grandes o los de curvas elípticas, que son más fuertes con llaves más
    chicas.
\end{itemize}

Los grupos que hoy se consideran seguros son:

\begin{itemize}
    \item \textbf{Grupo 14} (2048 bits): es el mínimo aceptable hoy.
    \item \textbf{Grupos 15 y 16} (3072 y 4096 bits): más fuertes.
    \item \textbf{Grupos 19, 20 y 21} (curvas elípticas de 256, 384 y 521 bits):
    los más recomendados ahorita. El 20 es el que pide la NSA en su suite CNSA.
\end{itemize}

En la práctica funcionó el grupo 5 porque es una simulación, pero en una red
real habría que usar mínimo el 14, y de preferencia el 19 o el 20.
\end{respuesta}
